% auto-generated by src/eval/make_tables.py -- do not edit by hand
\begin{table*}[t]
\centering
\small
\begin{tabular}{lrrrr}
\toprule
Predictor of followability & Pearson $r$ & 95\% CI & Williams $p$ & System $\rho$ \\
\midrule
IFS (source-anchored) & 0.13 & [-0.04, 0.29] & -- & 0.20 \\
Learned estimator & 0.53 & [0.37, 0.66] & $<$0.001 & 1.00 \\
\quad (Gemini outputs held out) & 0.55 & [0.39, 0.68] & $<$0.001 & 1.00 \\
\bottomrule
\end{tabular}
\caption{A learned per-dimension estimator recovers what IFS loses. The estimator is an LLM judge that scores the same four procedural dimensions (step, action, entity, quantity) yes/no from the source and translation; its composite is the mean of the four. Item-level Pearson $r$ is against human followability ($n{=}160$; CIs are 5{,}000-sample percentile bootstraps); Williams $p$ tests the gap to IFS on the shared human variable; System $\rho$ is the descriptive Spearman of each predictor's per-system means against the human ranking ($n{=}4$ systems, reported descriptively and without a $p$-value). The estimator predicts followability where IFS cannot ($r{=}0.53$ vs.\ 0.13) and recovers the human ordering ($\rho{=}1.00$), which saturated IFS does not ($\rho{=}0.20$). Holding out the judge model's own translations ($n{=}120$) leaves the recovery intact, so it is not self-preference.}
\label{tab:estimator}
\end{table*}
